\documentclass[11pt,border=4pt]{standalone}
% Estilo comum dos esquemáticos do relatório (circuitikz).
% Uso: \documentclass[11pt,border=4pt]{standalone} e, em seguida, % Estilo comum dos esquemáticos do relatório (circuitikz).
% Uso: \documentclass[11pt,border=4pt]{standalone} e, em seguida, % Estilo comum dos esquemáticos do relatório (circuitikz).
% Uso: \documentclass[11pt,border=4pt]{standalone} e, em seguida, \input{estilo_circuitos}.
\usepackage[T1]{fontenc}
\usepackage[utf8]{inputenc}
\usepackage[brazilian]{babel}
\usepackage{lmodern}
\usepackage{amsmath}
\usepackage[american]{circuitikz}
\usetikzlibrary{calc,arrows.meta}

% Cores das figuras originais
\definecolor{anot}{RGB}{0,135,200}      % anotações: valores calculados, setas, resistências vistas
\definecolor{sombra}{RGB}{220,242,253}  % região do transistor nos modelos T e pi
\definecolor{term}{RGB}{62,124,148}     % nomes de terminais destacados (B, C, E)

% Fios com 0,4 pt e corpo dos componentes com 0,8 pt, como nas figuras originais
\ctikzset{bipoles/thickness=2}
.
\usepackage[T1]{fontenc}
\usepackage[utf8]{inputenc}
\usepackage[brazilian]{babel}
\usepackage{lmodern}
\usepackage{amsmath}
\usepackage[american]{circuitikz}
\usetikzlibrary{calc,arrows.meta}

% Cores das figuras originais
\definecolor{anot}{RGB}{0,135,200}      % anotações: valores calculados, setas, resistências vistas
\definecolor{sombra}{RGB}{220,242,253}  % região do transistor nos modelos T e pi
\definecolor{term}{RGB}{62,124,148}     % nomes de terminais destacados (B, C, E)

% Fios com 0,4 pt e corpo dos componentes com 0,8 pt, como nas figuras originais
\ctikzset{bipoles/thickness=2}
.
\usepackage[T1]{fontenc}
\usepackage[utf8]{inputenc}
\usepackage[brazilian]{babel}
\usepackage{lmodern}
\usepackage{amsmath}
\usepackage[american]{circuitikz}
\usetikzlibrary{calc,arrows.meta}

% Cores das figuras originais
\definecolor{anot}{RGB}{0,135,200}      % anotações: valores calculados, setas, resistências vistas
\definecolor{sombra}{RGB}{220,242,253}  % região do transistor nos modelos T e pi
\definecolor{term}{RGB}{62,124,148}     % nomes de terminais destacados (B, C, E)

% Fios com 0,4 pt e corpo dos componentes com 0,8 pt, como nas figuras originais
\ctikzset{bipoles/thickness=2}

% Modelo genérico de amplificador de tensão dos dois estágios da Questão 1.
% Compilar com pdflatex; o PDF resultante vai para figuras/fig_c_q1_generico.pdf.
%
% Rótulos como nós explícitos (inner sep=0), postos a uma folga fixa do corpo de cada
% componente. Meias larguras dos símbolos do circuitikz (bipoles/length = 1,4 cm):
% resistor 0,21 cm, fonte controlada (losango) 0,49 cm, fonte independente (círculo) 0,42 cm.
\newcommand{\folgaLado}{3.85bp}   % do corpo do componente até o rótulo ao lado
\newcommand{\folgaPonto}{7.5bp}   % do centro do ponto de junção até a base do rótulo acima
\newcommand{\folgaRo}{5.95bp}     % do topo do zigue-zague de R_o até a base do rótulo
\begin{document}
\begin{circuitikz}
  % fonte de entrada e estágio 1
  \draw (0,0) to[V, invert] (0,2.6) -- (2.2,2.6);
  \draw (2.2,2.6) node[circ]{} to[R] (2.2,0);
  \draw (4.8,0) to[cV, invert] (4.8,2.6) to[R] (7.0,2.6) -- (8.4,2.6);
  % estágio 2 e carga
  \draw (8.4,2.6) node[circ]{} to[R] (8.4,0);
  \draw (11.0,0) to[cV, invert] (11.0,2.6) to[R] (13.2,2.6) -- (14.4,2.6);
  \draw (14.4,2.6) node[circ]{} to[R] (14.4,0);
  \draw (0,0) -- (14.4,0);
  \draw (0,0) node[ground]{};

  % rótulos dos componentes
  \node[anchor=mid east, inner sep=0pt] at ($(-0.42,1.3)+(-\folgaLado,0)$) {$v_i$};
  \node[anchor=mid west, inner sep=0pt] at ($(2.41,1.3)+(\folgaLado,0)$) {$R_{\mathrm{in}1}$};
  \node[anchor=mid west, inner sep=0pt] at ($(5.29,1.3)+(\folgaLado,0)$) {$A_{vo1}v_{i1}$};
  \node[anchor=south, inner sep=0pt] at ($(5.9,2.81)+(0,\folgaRo)$) {$R_{o1}$};
  \node[anchor=mid west, inner sep=0pt] at ($(8.61,1.3)+(\folgaLado,0)$) {$R_{\mathrm{in}2}$};
  \node[anchor=mid west, inner sep=0pt] at ($(11.49,1.3)+(\folgaLado,0)$) {$A_{vo2}v_{i2}$};
  \node[anchor=south, inner sep=0pt] at ($(12.1,2.81)+(0,\folgaRo)$) {$R_{o2}$};
  \node[anchor=mid west, inner sep=0pt] at ($(14.61,1.3)+(\folgaLado,0)$) {$R_L$};

  % tensões de nó em relação à terra
  \node[anchor=south, inner sep=0pt] at ($(2.2,2.6)+(0,\folgaPonto)$) {$v_{i1}$};
  \node[anchor=south, inner sep=0pt] at ($(8.4,2.6)+(0,\folgaPonto)$) {$v_{i2}$};
  \node[anchor=south, inner sep=0pt] at ($(14.4,2.6)+(0,\folgaPonto)$) {$v_o$};

  % limites dos estágios
  \draw[dashed, gray] (1.6,-0.45) rectangle (7.55,3.55);
  \draw[dashed, gray] (7.85,-0.45) rectangle (13.65,3.55);
  \node[gray, font=\small] at (4.575,3.93) {estágio 1 ($M_1$, dreno-comum)};
  \node[gray, font=\small] at (10.75,3.93) {estágio 2 ($M_2$, porta-comum)};
\end{circuitikz}
\end{document}
